\scanpage{15}

\margin{\deleted{\unsure{On le fait} (\unsure{fini}) \unsure{plus bas pour} $F$}}

\uline{Prop.} \deleted{(\unsure{1})} Soit $F\subset\mathfrak P(X)$ satisfaisant les conditions a) b) c) \inserted{d) e)} du \S\ 1, et soit $\Phi\subset F$ une partie finie de $F$. Alors il existe une décomposition cellulaire $F$-admissible de $X$, telle que les $Y\in\Phi$ sont des réunions de feuilles. Parmi \unsure{ttes ces|les ces} décompositions cellulaires, il y en a une plus \deleted{fine} grossière \deleted{que les autres}\inserted{(qui est donc \unsure{uniqt déterminée|unique déterminée})}.

\deleted{\unsure{Supp} $\Phi$ \unsure{stable par} $\unsure{f_0|b}$} Soit
\[
\unsure{X'|Z'}=\bigcup_{Y\in\Phi}\dot Y,
\]
c'est une réunion finie de parties fermées rares de $X$, donc est une partie \uline{rare} de $X$, soit $\deleted{\unsure{X'=X_2}}$ $U=\deleted{\unsure{X}}\complement X'$. Les traces \inserted{$Y\cap U$} sur $U$ des $Y\in\Phi$ sont des \deleted{ouverts} parties à la fois ouvertes et fermées \deleted{donc sont} en \uline{nb fini}. Il y a donc (lemme \inserted{ensembliste} immédiat) dans $U$ une plus \deleted{fine}\inserted{grossière} partition \inserted{(finie)} par des $U_i$ \deleted{ouverts} telles que les $Y\cap U$ soient des réunions \inserted{(finies)} de $U_i$; de plus, chaque $U_i$ est une intersection \inserted{(finie)} des tels ensembles $Y\cap U$, \inserted{donc est aussi ouvert et fermé.} Ceci étant, si la prop. est prouvée pour \deleted{\unsure{les}} $(X,\Phi)$ remplacé par
\[
\left(X',\Phi'=\{Y\cap X'\mid Y\in\Phi\}\inserted{\cup\{\overline U_i\cap X'\}_{i\in\unsure{I}}}\right),
\]

\mspage{16}
